\begin{solution}{4.01}
(1) $\delta$ should be the sum of the angles of deviation produced in the two refractions. From
Fig.~4.01-1,
\begin{equation*}
\delta = i-r+r'-i'.
\end{equation*}
From the geometry,
\begin{equation}
r+i' = A. \tag{1}
\end{equation}
Hence
\begin{equation}
\delta = i+r'-A. \tag{2}
\end{equation}
By the law of refraction,
\begin{align}
\sin i &= n\sin r, \tag{3}\\
\sin r' &= n\sin i'. \tag{4}
\end{align}

\begin{figbox}{4.01-1}
\begin{tikzpicture}[scale=1.15]
  % prism
  \draw[thickline] (0,2.2) -- (-1.6,0) -- (1.6,0) -- cycle;
  \node[above] at (0,2.26) {$A$};
  % normals
  \draw[thinline,dashed] (-1.005,1.753) -- (-0.115,1.107);
  \draw[thinline,dashed] (0.159,1.046) -- (1.049,1.693);
  % incident ray
  \draw[thickline,->] (-2.90,0.647) -- (-0.56,1.43);
  \draw[dashedline] (-0.56,1.43) -- (0.83,1.891);
  % ray inside the prism
  \draw[thickline,->] (-0.56,1.43) -- (0.604,1.3695);
  % emergent ray
  \draw[dashedline] (0.604,1.3695) -- (0.167,1.672);
  \draw[thickline,->] (0.604,1.3695) -- (1.95,0.436);
  % angle i at the first face
  \draw ( -0.56,1.43) ++(144:0.42) arc (144:198.4:0.42);
  \node at (-1.26,1.60) {$i$};
  % angle r at the first face
  \draw (-0.56,1.43) ++(-36:0.46) arc (-36:-3:0.46);
  \node at (-0.20,1.13) {$r$};
  % angle i' at the second face
  \draw (0.604,1.3695) ++(177:0.42) arc (177:216:0.42);
  \node at (0.42,1.05) {$i'$};
  % angle r' at the second face
  \draw (0.604,1.3695) ++(36:0.62) arc (36:-34.7:0.62);
  \node at (1.52,1.70) {$r'$};
  % deviation
  \draw (0.167,1.672) ++(18.4:0.46) arc (18.4:-34.7:0.46);
  \node at (0.86,1.90) {$\delta$};
\end{tikzpicture}
\end{figbox}

From (1), $i'=A-r$; substituting this into (4) and using (3), we have
\begin{align*}
\sin r' &= n\sin(A-r) = n\left(\sin A\cos r-\cos A\sin r\right)\\
&= \sin A\sqrt{n^2-\sin^2 i}-\cos A\sin i.
\end{align*}
Substituting this into (2), we obtain
\begin{equation*}
\delta = i+\arcsin\left(\sin A\sqrt{n^2-\sin^2 i}-\cos A\sin i\right)-A.
\end{equation*}

(2) For $\delta$ to be minimum we must have $\dd\delta/\dd i=0$. Differentiating (2) and setting the
derivative equal to zero,
\begin{equation*}
\frac{\dd\delta}{\dd i} = 1+\frac{\dd r'}{\dd i} = 0,
\end{equation*}
we obtain
\begin{equation}
\frac{\dd r'}{\dd i} = -1. \tag{5}
\end{equation}

Differentiating (3) and (4),
\begin{align*}
\cos i\,\dd i &= n\cos r\,\dd r,\\
\cos r'\,\dd r' &= n\cos i'\,\dd i'.
\end{align*}
Dividing the first of these by the second, and using the result $\dd r=-\dd i'$ obtained by
differentiating (1) together with (5), we obtain
\begin{equation*}
\frac{\cos i}{\cos r} = \frac{\cos r'}{\cos i'}.
\end{equation*}
Substituting (3) and (4) into this, we obtain
\begin{equation*}
\frac{1-\sin^2 i}{1-\sin^2 r'} = \frac{n^2-\sin^2 i}{n^2-\sin^2 r'}.
\end{equation*}
Since $n\neq 1$, for the above equation to hold we must have
\begin{equation*}
i = r',
\end{equation*}
and hence
\begin{equation*}
r = i'.
\end{equation*}
It is thus seen that when the incident and emergent rays are symmetric, the deviation $\delta$
takes an extremal value $\delta_m$. Further verification shows that $\delta$ then takes its
minimum value, i.e. $\delta_m=\delta_{\min}$.

In the case $\delta=\delta_m$, from (1) and (2),
\begin{align*}
r &= i' = A/2,\\
i &= r' = (\delta_m+A)/2.
\end{align*}
Substituting into (3), we obtain
\begin{equation}
n = \frac{\sin\left[(\delta_m+A)/2\right]}{\sin(A/2)}, \tag{6}
\end{equation}
or
\begin{equation*}
\delta_m = 2\arcsin\left(n\sin\frac{A}{2}\right)-A.
\end{equation*}

\begin{commentary}
In experiments the minimum angle of deviation can be measured precisely with ease, so equation
(6) is often used to determine the refractive index of the material of a prism.

If it is only required to prove that the deviation takes an extremal value when the light path is
symmetric, a simpler method can also be used; it is described briefly below.

Let the symmetric path be $1$--$1'$, as shown in Fig.~4.01-2. For this path the angle of
incidence and the angle of emergence are equal ($i=r'$); let its deviation be $\delta$. Now let
the light be incident at a slightly larger angle of incidence; then its angle of emergence is
slightly smaller. Denote this path by $2$--$2'$ and its deviation by $\delta'\neq\delta$. By the
principle of reversibility of the light path, the light can also travel in the reverse direction
along the same path $2'$--$2$. This is equivalent to reducing the angle of incidence of the light
relative to its value for the symmetric path (and increasing the angle of emergence), while the
deviation of $2'$--$2$ remains $\delta'$. Since the deviation is the same whether the angle of
incidence is increased or decreased relative to its value for the symmetric path, it is evident
that the deviation of the symmetric path must be extremal.

\begin{figbox}{4.01-2}
\begin{tikzpicture}[scale=1.15]
  \draw[thickline] (0,2.5) -- (-1.8,0) -- (1.8,0) -- cycle;
  \node[above] at (0,2.52) {$A$};
  % symmetric path 1-1'
  \draw[thickline,->] (-3.10,0.91) -- (-0.792,1.40);
  \draw[thickline,->] (-0.792,1.40) -- (0.792,1.40);
  \draw[thickline,->] (0.792,1.40) -- (3.10,0.91);
  % perturbed path 2-2'
  \draw[thickline,dashed,->] (-3.10,0.60) -- (-0.922,1.22);
  \draw[thickline,dashed,->] (-0.922,1.22) -- (0.93,1.16);
  \draw[thickline,dashed,->] (0.93,1.16) -- (3.10,0.66);
  % labels
  \node[left] at (-3.16,0.56) {$2$};
  \node[left] at (-3.16,0.88) {$1$};
  \node[right] at (3.12,0.62) {$2'$};
  \node[right] at (3.12,0.88) {$1'$};
  % deviation
  \draw[dashedline] (0.792,1.40) -- (-1.30,1.86);
  \draw[dashedline] (0.93,1.16) -- (-1.15,1.90);
  \node at (0.34,2.02) {$\delta$};
\end{tikzpicture}
\end{figbox}
\end{commentary}
\end{solution}

\begin{solution}{4.02}
(1) From symmetry it is known that a ray of light travelling in an isosceles prism must be parallel
to the base, as shown in Fig.~4.02-1. Applying the law of refraction at the first and the second
interface, we have
\begin{align}
\sin i_1 &= n_1\sin r_1, \tag{1}\\
n_1\sin i_2 &= n_2\sin r_2. \tag{2}
\end{align}
From the geometry,
\begin{align*}
i_1 &= \frac{\pi}{2}-\frac{\alpha}{2},\\
i_2 &= \frac{\pi}{2}-r_1,\\
r_2 &= \frac{\alpha}{2}.
\end{align*}
Substituting these into (1) and (2), we obtain
\begin{align*}
\cos\frac{\alpha}{2} &= n_1\sin r_1,\\
n_1\cos r_1 &= n_2\sin\frac{\alpha}{2}.
\end{align*}
From these two equations we can solve for
\begin{equation*}
\sin\frac{\alpha}{2} = \sqrt{\frac{n_1^2-1}{n_2^2-1}},
\end{equation*}
that is,
\begin{equation}
\alpha = 2\arcsin\sqrt{\frac{n_1^2-1}{n_2^2-1}}. \tag{3}
\end{equation}

\begin{figbox}{4.02-1}
\begin{tikzpicture}[scale=1.1]
  \draw[thickline] (0,3.0) -- (-2.0,0) -- (2.0,0) -- cycle;
  \draw[dashedline] (0,3.0) -- (0,0);
  \node[above] at (0,3.04) {$A$};
  % apex angles
  \node[scale=0.9] at (0.34,2.80) {$\alpha$};
  \draw (0,3.0) ++(236.31:0.62) arc (236.31:270:0.62);
  \node[scale=0.8] at (-0.21,2.39) {$\frac{\alpha}{2}$};
  % incident ray
  \draw[thickline,->] (-3.40,0.85) -- (-0.9667,1.55);
  % normals
  \draw[thinline,dashed] (-0.9667,1.55) ++(146.31:0.75) -- ++(326.31:1.40);
  \draw[thinline,dashed] (0.9667,1.55) ++(33.69:0.85) -- ++(213.69:1.55);
  % ray inside the prism
  \draw[thickline,->] (-0.9667,1.55) -- (0.9667,1.55);
  % emergent ray
  \draw[thickline,->] (0.9667,1.55) -- (2.72,0.82);
  % angles
  \draw (-0.9667,1.55) ++(146.31:0.60) arc (146.31:196.05:0.60);
  \node at (-1.70,1.72) {$i_1$};
  \draw (-0.9667,1.55) ++(-33.69:0.50) arc (-33.69:0:0.50);
  \node at (-0.52,1.22) {$r_1$};
  \draw (0.9667,1.55) ++(180:0.50) arc (180:213.69:0.50);
  \node at (0.52,1.20) {$i_2$};
  \draw (0.9667,1.55) ++(33.69:0.80) arc (33.69:-22.6:0.80);
  \node at (1.94,1.64) {$r_2$};
\end{tikzpicture}
\end{figbox}

(2) From the table, for the $F$ line,
\begin{equation*}
n_1^F = 1.515\,65,\qquad n_2^F = 1.774\,75.
\end{equation*}
Substituting into (3), we obtain
\begin{equation*}
\alpha = 2\arcsin\sqrt{\frac{1.515\,65^2-1}{1.774\,75^2-1}} = 101.94\degree.
\end{equation*}
For the $C$ line and the $G$ line,
\begin{align*}
n_1^C &= 1.507\,60,\qquad n_2^C = 1.747\,32,\\
n_1^G &= 1.520\,16,\qquad n_2^G = 1.792\,05.
\end{align*}
The angles of refraction are calculated one after another and listed in the table below. The last
column gives the deviation $\delta$ of the emergent light with respect to the incident light, taken
positive in the anticlockwise sense.
\end{solution}

\begin{solution}{4.11}
(1) Consider an arbitrary ray $PA$ issuing from $P$, and let the angle it makes with the principal
axis be $\theta$, as shown in Fig.~4.11-1. Let the angle of incidence of the ray at the spherical
surface be $i$ and the angle of refraction be $r$. By the law of refraction,
\begin{equation}
\frac{\sin i}{\sin r} = \frac{1}{n}. \tag{1}
\end{equation}

\begin{figbox}{4.11-1}
\begin{tikzpicture}[scale=0.95]
  \draw[thickline] (-2.9,0) -- (2.1,0);
  \draw[thickline] (0,0) circle (1.6);
  \fill (0,0) circle (1.3pt) node[below] {$C$};
  \fill (-1.0667,0) circle (1.3pt) node[below] {$P$};
  \fill (-2.4,0) circle (1.2pt) node[below] {$Q$};
  % incident ray inside the sphere
  \draw[thickline,->] (-1.0667,0) -- (0.666,1.455);
  % refracted ray
  \draw[thickline,->] (0.666,1.455) -- (1.90,2.043);
  % backward extension of the refracted ray
  \draw[dashedline] (0.666,1.455) -- (-2.40,0);
  % radius (normal) at A
  \draw[dashedline] (0.666,1.455) -- (0,0);
  \fill (0.666,1.455) circle (1.3pt) node[above left,inner sep=1pt] {$A$};
  % theta at P
  \draw (-1.0667,0) ++(0:0.50) arc (0:40:0.50);
  \node at (-0.50,0.22) {$\theta$};
  % i and r at A
  \draw (0.666,1.455) ++(220:0.45) arc (220:245.4:0.45);
  \node at (0.28,0.96) {$i$};
  \draw (0.666,1.455) ++(65.4:0.62) arc (65.4:25.4:0.62);
  \node at (1.28,2.05) {$r$};
  % media
  \node at (0.32,0.52) {$n$};
  \node at (1.18,1.70) {$1$};
  % dimensions
  \draw[dim] (0,-0.30) -- (1.6,-0.30) node[midway,below] {$R$};
  \draw[dim] (-2.4,-0.75) -- (0,-0.75) node[midway,below] {$nR$};
\end{tikzpicture}
\end{figbox}

In $\triangle PCA$, applying the sine rule, we have
\begin{equation*}
\frac{\sin i}{R/n} = \frac{\sin\theta}{R},
\end{equation*}
or
\begin{equation}
\frac{\sin i}{\sin\theta} = \frac{1}{n}. \tag{2}
\end{equation}

Comparing (1) and (2), in the case where $\theta$ is an acute angle ($r$ must be acute), we have
\begin{equation*}
\theta = r.
\end{equation*}
But $r=\angle QAC$, where $Q$ is the point of intersection of the backward extension of the
refracted ray with the principal axis; hence $\triangle QAC\sim\triangle APC$, and
$\overline{QC}/R=R/(R/n)=n$, i.e.
\begin{equation*}
\overline{QC} = nR,
\end{equation*}
independently of the angle $\theta$, which was to be proved.

(2) The geometrical configuration of the two lenses is shown in Fig.~4.11-2.

\begin{figbox}{4.11-2}
\begin{tikzpicture}[scale=0.95]
  \draw[thickline] (-6.6,0) -- (2.9,0);
  % objective lens L1
  \draw[thickline] (-3.00,0) circle (0.60);
  \node at (-2.78,0.16) {$n_1$};
  \draw[thinline,dashed] (-3.00,0) -- (-3.00,0.60);
  \node at (-3.42,0.56) {$R_1$};
  \fill (-3.00,0) circle (1.2pt) node[below] {$C_1$};
  \node[above] at (-3.00,0.66) {$L_1$};
  \fill (-2.40,0) circle (1.0pt) node[above right,inner sep=1pt] {$O_1$};
  % second lens L2
  \draw[thickline] (-4.60,0) ++(10.8:4.00) arc (10.8:-10.8:4.00);
  \draw[thickline] (-0.90,0) ++(56.4:0.90) arc (56.4:-56.4:0.90);
  \draw[thickline] (-0.67,0.75) -- (-0.40,0.75);
  \draw[thickline] (-0.67,-0.75) -- (-0.40,-0.75);
  \node[above] at (-1.95,0.82) {$L_2$};
  \node at (-0.36,0.60) {$n_2$};
  \fill (-0.90,0) circle (1.2pt) node[below,inner sep=2pt] {$C_3$};
  \fill (-0.60,0) circle (1.0pt) node[below right,inner sep=1pt] {$O_2$};
  \fill (0,0) circle (1.0pt) node[below right,inner sep=1pt] {$O_3$};
  % object and images
  \fill (-6.00,0) circle (1.3pt) node[below] {$S$};
  \fill (-4.60,0) circle (1.3pt) node[below] {$S'$};
  \fill (2.40,0) circle (1.3pt) node[below] {$S''$};
  % aperture angle at S
  \draw[thickline] (-6.00,0) -- (-3.58,0.42);
  \draw[thickline] (-6.00,0) -- (-3.58,-0.42);
  \draw (-6.00,0) ++(0:1.10) arc (0:10.4:1.10);
  \node at (-4.92,0.42) {$\alpha$};
  % dimensions
  \draw[dim] (-2.40,-0.95) -- (-0.60,-0.95) node[midway,below] {$d_1$};
  \draw[dim] (-0.60,-1.30) -- (0,-1.30) node[midway,below] {$d_2$};
  \node at (-1.30,0.72) {$R_2$};
  \node at (-0.12,0.95) {$R_3$};
\end{tikzpicture}
\end{figbox}

(i) Let $C_1$ be the centre of the sphere of the objective lens ($L_1$). The light emitted from $S$
(the object) is refracted by $L_1$ and forms an image at $S'$. From part (1) of this problem,
$\overline{S'C_1}=n_1R_1$, so
\begin{equation*}
\overline{S'O_1} = n_1R_1+R_1 = (n_1+1)R_1.
\end{equation*}
The light emitted from $S'$ should enter the front spherical surface of the second lens ($L_2$)
without refraction. Hence $S'$ is the centre of the front spherical surface, and the radius of the
front spherical surface is
\begin{equation*}
R_2 = \overline{S'O_2} = \overline{S'O_1}+d_1 = (n_1+1)R_1+d_1.
\end{equation*}
In order that $S'$ be located at the aplanatic point of $L_2$, and that the central thickness of
$L_2$ be $d_2$, we must have
\begin{equation*}
\overline{S'O_3} = R_3+\frac{R_3}{n_2}
\end{equation*}
and
\begin{equation*}
\overline{S'O_3} = R_2+d_2 = (n_1+1)R_1+d_1+d_2.
\end{equation*}
From these two equations we can solve for the radius of the back spherical surface:
\begin{equation*}
R_3 = \frac{(n_1+1)R_1+d_1+d_2}{1+\dfrac{1}{n_2}}.
\end{equation*}
In Fig.~4.11-2, $S''$ is the image of $S'$ formed by $L_2$ [$\overline{S''O_3}=(n_2+1)R_3$], and
$C_3$ is the centre of the back spherical surface.

(ii) The aperture angle $\alpha$ of the light emitted from $S$ satisfies
\begin{equation*}
\tan\alpha = \frac{R_1}{R_1/n_1} = n_1.
\end{equation*}
Let the aperture angle of the image $S'$ formed by $S$ through $L_1$ be $\beta_1$; then
\begin{equation*}
\tan\beta_1 = \frac{R_1}{n_1R_1} = \frac{1}{n_1}.
\end{equation*}
$\beta_1$ is also the aperture angle of the object point of $L_2$, and $\beta$ is the aperture angle
of the final image $S''$. From Fig.~4.11-2 it is easily seen that
\begin{equation*}
\tan\beta = \frac{\overline{S'A}\sin\beta_1}{\overline{S'A}\cos\beta_1+\overline{S'S''}}
\approx \frac{R_2\sin\beta_1}{R_2\cos\beta_1+\overline{S'S''}}
\qquad\text{(the edge thickness of the lens being neglected)}.
\end{equation*}
Moreover,
\begin{equation*}
\overline{S'S''} = n_2R_3-\frac{1}{n_2}R_3,
\end{equation*}
therefore
\begin{equation*}
\tan\beta = \frac{R_2\sin\beta_1}{R_2\cos\beta_1+\left(n_2-\dfrac{1}{n_2}\right)R_3}.
\end{equation*}

(iii) $k = k_1\cdot k_2 = \dfrac{n_1R_1}{R_1/n_1}\cdot\dfrac{n_2R_2}{R_2/n_2} = n_1^2\cdot n_2^2$.

(iv) Substituting the data, we obtain $R_2=9.5\unitv{mm}$, $R_3=6.77\unitv{mm}$,
$\beta=20.0\degree$, $k=5.76$.

\begin{commentary}
This problem is an application of the concept of aplanatic points to the microscope objective. Such
an objective is called an oil-immersion objective and is one of the main types of microscope
objective.
\end{commentary}
\end{solution}

\begin{solution}{4.12}
(1) From Fig.~4.12-1, the deviation of the light emitted from $S$ after two refractions by the
prism is
\begin{equation}
\delta = i-r+r'-i' = i+r'-(r+i'). \tag{1}
\end{equation}
But
\begin{equation}
r+i' = \alpha. \tag{2}
\end{equation}
From the statement of the problem, $i$ is a small angle, hence $r$, $i'$ and $r'$ are also small
angles; therefore, by the law of refraction,
\begin{equation*}
i = nr,\qquad r' = ni'.
\end{equation*}
Substituting these into (1) and using (2), we obtain
\begin{align}
\delta &= nr+ni'-(r+i')\notag\\
&= (n-1)\alpha. \tag{3}
\end{align}

\begin{figbox}{4.12-1}
\begin{tikzpicture}[scale=1.45]
  \draw[thickline] (0,3.0) -- (0,-0.8) -- (0.85,-0.8) -- cycle;
  \node[left] at (0,3.0) {$A$};
  \draw (0,3.0) ++(270:0.55) arc (270:257.4:0.55);
  \node at (0.30,2.62) {$\alpha$};
  % normals
  \draw[thinline,dashed] (-0.55,1.35) -- (0.55,1.35);
  \draw[thinline,dashed] (0.334,1.509) ++(12.6:0.55) -- ++(192.6:1.10);
  % incident ray
  \draw[thickline,->] (-1.60,0.00) -- (0,1.35);
  \draw[dashedline] (0,1.35) -- (0.48,1.75);
  % ray inside the prism
  \draw[thickline,->] (0,1.35) -- (0.334,1.509);
  % emergent ray
  \draw[thickline,->] (0.334,1.509) -- (1.95,1.315);
  % angles at the first face
  \draw (0,1.35) ++(180:0.42) arc (180:220.2:0.42);
  \node at (-0.72,1.16) {$i$};
  \draw (0,1.35) ++(0:0.55) arc (0:25.4:0.55);
  \node at (0.30,1.20) {$r$};
  % angles at the second face
  \draw (0.334,1.509) ++(192.6:0.34) arc (192.6:205.4:0.34);
  \node at (0.02,1.28) {$i'$};
  \draw (0.334,1.509) ++(12.6:0.62) arc (12.6:-6.8:0.62);
  \node at (1.02,1.18) {$r'$};
  % deviation
  \node at (0.66,1.74) {$\delta$};
\end{tikzpicture}
\end{figbox}

Since $\delta$ is independent of the angle of incidence, the backward extensions $S_1A$ and $S_1B$ of
any two rays $SA$ and $SB$ emitted from $S$ with an angle $\Delta\theta$ between them, after
refraction through the upper half of the prism, still make the angle $\Delta\theta$ (see
Fig.~4.12-2), and the arc which they subtend is the same arc $AB$; hence the distance of $S_1$
(i.e. the image of $S$) from the prism is also $a$. In the same way, the image $S_2$ formed by $S$
through refraction at the lower half of the prism is also at a distance $a$ from the prism. Thus
(using (3)) the distance between the two coherent sources is
\begin{equation*}
\overline{S_1S_2} = 2a\delta = 2a(n-1)\alpha.
\end{equation*}

\begin{figbox}{4.12-2}
\begin{tikzpicture}[scale=1.1]
  % biprism
  \draw[thickline] (0,1.4) -- (-0.28,0) -- (0.28,0) -- cycle;
  \draw[thickline] (0,-1.4) -- (-0.28,0) -- (0.28,0) -- cycle;
  % screen
  \draw[thickline] (3.6,-1.5) -- (3.6,1.5);
  % axis
  \draw[dashedline] (-4.0,0) -- (4.0,0);
  % object and virtual images
  \fill (-2.90,0) circle (1.3pt) node[below] {$S$};
  \fill (-3.65,0.50) circle (1.2pt) node[above left,inner sep=1pt] {$S_1$};
  \fill (-3.65,-0.50) circle (1.2pt) node[below left,inner sep=1pt] {$S_2$};
  % rays from S through the upper half prism
  \draw[thickline] (-2.90,0) -- (-0.15,0.62);
  \draw[thickline] (-0.15,0.62) -- (3.6,0.30);
  \draw[thickline] (-2.90,0) -- (-0.15,0.12);
  \draw[thickline] (-0.15,0.12) -- (3.6,0.30);
  \fill (-0.15,0.62) circle (1.1pt) node[above right,inner sep=1pt] {$A$};
  \fill (-0.15,0.12) circle (1.1pt) node[below right,inner sep=1pt] {$B$};
  % rays from S through the lower half prism
  \draw[thickline] (-2.90,0) -- (-0.15,-0.62);
  \draw[thickline] (-0.15,-0.62) -- (3.6,0.30);
  \draw[thickline] (-2.90,0) -- (-0.15,-0.12);
  \draw[thickline] (-0.15,-0.12) -- (3.6,0.30);
  % backward extensions
  \draw[dashedline] (-0.15,0.62) -- (-3.65,0.50);
  \draw[dashedline] (-0.15,0.12) -- (-3.65,0.50);
  \draw[dashedline] (-0.15,-0.62) -- (-3.65,-0.50);
  \draw[dashedline] (-0.15,-0.12) -- (-3.65,-0.50);
  % observation point
  \fill (3.6,0.30) circle (1.3pt) node[above right,inner sep=1pt] {$P$};
  \fill (3.6,0) circle (1.2pt) node[below right,inner sep=1pt] {$O$};
  \draw[dim] (3.42,0) -- (3.42,0.30) node[midway,left] {$x$};
  % aperture angle
  \draw (-2.90,0) ++(2.5:1.20) arc (2.5:12.7:1.20);
  \node at (-2.05,0.58) {$\Delta\theta$};
  % distances
  \draw[dim] (-2.90,-1.75) -- (0,-1.75) node[midway,below] {$a$};
  \draw[dim] (0,-1.75) -- (3.6,-1.75) node[midway,below] {$b$};
\end{tikzpicture}
\end{figbox}

It is not difficult to see that the optical path difference from $S_1$ and $S_2$ to the point $P$
(the distance of $P$ from the central point $O$ being $x$) is
\begin{equation}
\Delta = \overline{S_1S_2}\cdot\frac{x}{a+b} = \frac{2a(n-1)\alpha x}{a+b}. \tag{4}
\end{equation}
Setting $\Delta=m\lambda$, we obtain the distance $x$ of the $m$-th bright fringe:
\begin{equation*}
x = \frac{m\lambda(a+b)}{2a(n-1)\alpha}.
\end{equation*}

(2) When $S$ is moved up to $S'$, imagine that a thin plate of thickness $d$ is cut away from the
biprism along the dashed line; then the zero-order fringe would move up to $O'$ (Fig.~4.12-3). The
presence of the thin plate of thickness $d$ produces for $S'$ an additional optical path difference
$(n-1)d$ in the lower half of the prism, so that the zero-order fringe moves down from $O'$ by a
distance $y'$, and [cf. (4)]
\begin{equation*}
\frac{2a(n-1)\alpha y'}{a+b} = (n-1)d.
\end{equation*}
From the geometrical relation $d=(h/\cos\alpha)\cdot\sin 2\alpha\approx 2h\alpha$; substituting this
into the above equation gives $y'=(a+b)h/a$. Hence the upward displacement of $S$ to $S'$ finally
makes the zero-order fringe move down from the point $O$ by a distance
\begin{equation*}
y = y'-h = \frac{b}{a}h.
\end{equation*}

\begin{figbox}{4.12-3}
\begin{tikzpicture}[scale=0.95]
  \draw[thickline] (0,1.3) -- (-0.28,0) -- (0.28,0) -- cycle;
  \draw[thickline] (0,-1.3) -- (-0.28,0) -- (0.28,0) -- cycle;
  \fill[black!30] (-0.30,-0.09) rectangle (0.30,0.09);
  \draw[thickline] (-0.30,-0.09) rectangle (0.30,0.09);
  \node at (-0.52,0.16) {$d$};
  \draw[thickline] (2.30,-1.0) -- (2.30,1.20);
  \draw[dashedline] (-3.20,0.45) -- (2.30,0.45);
  \draw[dashedline] (-3.20,0) -- (2.30,0);
  \fill (-2.55,0.45) circle (1.3pt) node[above] {$S'$};
  \fill (-2.55,0) circle (1.3pt) node[below] {$S$};
  \fill (2.30,0.45) circle (1.2pt) node[right] {$O'$};
  \fill (2.30,0) circle (1.2pt) node[right] {$O$};
  \draw[dim] (-2.25,0) -- (-2.25,0.45) node[midway,right] {$h$};
\end{tikzpicture}
\end{figbox}

\begin{commentary}
There are three keys to the solution of this problem. First, to prove that the light emitted from
the source $S$ is deviated through a small angle $(n-1)\alpha$ after refraction by the prism;
second, to show that the image of $S$ formed by refraction through the prism is still at a distance
$a$ from the prism, and thereby to obtain the distance between the two coherent sources; finally,
to find the distance through which the zero-order fringe moves down when the source moves up, by
imagining a thin plate cut away from the prism.
\end{commentary}
\end{solution}

\begin{solution}{4.13}
(1) From the conditions of the problem the geometrical light path is as shown in Fig. 4.13(a).
On the screen $M$ one obtains the interference fringes of the parallel beam 1, formed from the
light of $S$ after refraction by $L_1$, and of the divergent beam 2, emitted from the image
point $S'$ after refraction and imaging by $L_2$; interference occurs only in the region above
the optical axis, so the fringes are concentric semicircles centred on the optical axis, as
shown on the right of Fig. 4.13(a). At the centre $O$ the optical path difference of the two
beams is originally zero, but the converging beam 2 undergoes a phase jump of $\pi$ when it
passes through its focal point $S'$, so that the optical path difference of the two beams at
the centre $O$ is $\lambda/2$; the centre is therefore dark.

The occurrence of a $\pi$ phase jump when a converging beam passes through its focal point can
be explained as follows. For the divergent wave emitted from the point $S$ (Fig. 4.13(b)), in
order to obtain by Fresnel's half-period-zone method the correct phase of the wave disturbance
arriving at $P$, one must take the phase of every secondary wavelet emitted from the wavefront
$\Sigma$ to be ahead of the phase of the wave arriving at $\Sigma$ from $S$ by $\pi/2$. For the
converging wave $\Sigma'$ focused at $S'$, when the observation point $P$ lies between $\Sigma'$
and $S'$ (Fig. 4.13(c)(1)), it is enough to assume that the phase of the secondary wavelets on
$\Sigma'$ is ahead of the phase of the wave on $\Sigma'$ by $\pi/2$; the correct phase of the
wave disturbance at $P$ is still obtained, because the distances from the half-period zones on
$\Sigma'$ to $P$ still increase gradually. But when $P$ lies outside $\Sigma'$ and $S'$
(Fig. 4.13(c)(2)), since, counting from the first half-period zone, the distances from the
successive half-period zones to $P$ decrease one after another, their phases advance one after
another; adding the $\pi/2$ by which the wavelets on $\Sigma'$ were originally ahead, the phase
of the wave disturbance at $P$ will differ by $\pi$ from the phase obtained from $\Sigma'$
together with the optical path. This is the origin of the $\pi$ phase jump.

(2) The two beams have equal optical paths at $S'$. At a point $P$ of the screen $M$ at a
distance $x$ from $O$, beam 1 has in addition traversed the optical path $S'O$, while beam 2 has
traversed $S'P$. Taking further into account the $\pi$ phase jump of beam 2 at $S'$, the optical
path difference of the two beams at $P$ is
\begin{equation*}
\Delta = S'P + \frac{\lambda}{2} - S'O = \sqrt{f^2+x^2} + \frac{\lambda}{2} - f .
\end{equation*}
Putting $\Delta = m\lambda$, we obtain the position of the $m$-th bright fringe:
\begin{equation*}
x = \sqrt{\left(m-\frac{1}{2}\right)^2\lambda^2 + 2\left(m-\frac{1}{2}\right)\lambda f } .
\end{equation*}
It is seen that the interference fringes become denser towards the outside.

\begin{commentary}
The $\pi$ phase jump undergone by light on passing through a focus is easily overlooked, and
special attention should be paid to it.
\end{commentary}

\begin{figbox}{4.13(a)}
\begin{tikzpicture}[scale=1.0]
  % optical axis
  \draw[axis] (-0.2,0) -- (10.4,0);
  \fill (0.4,0) circle (1.2pt);
  \node[below left] at (0.4,0) {$S$};
  % lens L1 (above the axis)
  \draw[thickline] (3.8,1.80) .. controls (4.12,1.42) and (4.12,0.58) .. (3.8,0.20);
  \draw[thickline] (3.8,1.80) .. controls (3.48,1.42) and (3.48,0.58) .. (3.8,0.20);
  \node[above] at (3.8,1.80) {$L_1$};
  % lens L2 (below the axis)
  \draw[thickline] (6.2,-1.80) .. controls (6.52,-1.42) and (6.52,-0.58) .. (6.2,-0.20);
  \draw[thickline] (6.2,-1.80) .. controls (5.88,-1.42) and (5.88,-0.58) .. (6.2,-0.20);
  \node[below] at (6.2,-1.80) {$L_2$};
  % beam 1
  \draw[thickline] (0.4,0) -- (3.80,1.55);
  \draw[thickline] (3.80,1.55) -- (9.8,1.55);
  \node at (7.6,1.78) {1};
  % beam 2
  \draw[thickline] (0.4,0) -- (6.20,-1.55);
  \draw[thickline] (6.20,-1.55) -- (8.2,0);
  \fill[black!8] (8.2,0) -- (9.8,1.70) -- (9.8,-1.70) -- cycle;
  \draw[thickline] (8.2,0) -- (9.8,1.70);
  \draw[thickline] (8.2,0) -- (9.8,-1.70);
  \node at (9.25,-1.05) {2};
  \fill (8.2,0) circle (1.2pt);
  \node[below] at (8.2,-0.06) {$S'$};
  % screen
  \draw[thickline] (9.8,-2.0) -- (9.8,2.2);
  \node[above] at (9.8,2.2) {$M$};
  \fill (9.8,0) circle (1.1pt);
  \node[below left] at (9.8,-0.02) {$O$};
  \fill (9.8,2.0) circle (1.1pt);
  \node[right] at (9.8,2.0) {$A$};
  \fill (9.8,0.80) circle (1.1pt);
  \node[right] at (9.8,0.82) {$P$};
  % fringes on the screen
  \foreach \r in {0.22,0.38,0.54,0.70}
    \draw[thickline] (9.8,0) ++(-90:\r) arc (-90:90:\r);
  % distances
  \draw[dim] (0.4,-2.45) -- (3.8,-2.45) node[midway,below] {$f$};
  \draw[dim] (6.2,-2.45) -- (9.8,-2.45) node[midway,below] {$2f$};
\end{tikzpicture}
\end{figbox}

\begin{figbox}{4.13(b)}
\begin{tikzpicture}[scale=0.85]
  \fill (0,0) circle (1.2pt);
  \node[left] at (0,0) {$S$};
  \draw[thickline] (2.6,-1.60) arc (-31.6:31.6:3.0);
  \node[above] at (2.62,1.62) {$\Sigma$};
  \fill (3.0,0) circle (1.1pt);
  \node[left,fill=white,inner sep=1pt] at (2.92,0) {$A$};
  \draw[thickline] (0,0) -- (2.6,1.55);
  \draw[thickline] (0,0) -- (2.6,-1.55);
  \draw[thickline] (2.6,1.55) -- (5.5,0);
  \draw[thickline] (2.6,-1.55) -- (5.5,0);
  \fill (5.5,0) circle (1.1pt);
  \node[right] at (5.5,0) {$P$};
\end{tikzpicture}
\end{figbox}

\begin{figbox}{4.13(c)}
\begin{tikzpicture}[scale=0.8]
  % panel (1)
  \begin{scope}
    \fill (5.4,0) circle (1.1pt);
    \node[right,fill=white,inner sep=1pt] at (5.4,0) {$S'$};
    \draw[thickline] (5.4,0) ++(140:3.5) arc (140:220:3.5);
    \node[above] at (2.75,2.30) {$\Sigma'$};
    \fill (1.9,0) circle (1.1pt);
    \node[below left,fill=white,inner sep=1pt] at (1.9,0) {$A$};
    \fill (3.4,0) circle (1.1pt);
    \node[above] at (3.4,0.05) {$P$};
    \draw[thickline] (2.72,2.25) -- (5.4,0);
    \draw[thickline] (2.72,-2.25) -- (5.4,0);
    \node at (2.7,-2.6) {(1)};
  \end{scope}
  % panel (2)
  \begin{scope}[xshift=7.0cm]
    \fill (5.4,0) circle (1.1pt);
    \node[below,fill=white,inner sep=1pt] at (5.4,0) {$S'$};
    \draw[thickline] (5.4,0) ++(140:3.5) arc (140:220:3.5);
    \node[above] at (2.75,2.30) {$\Sigma'$};
    \fill (1.9,0) circle (1.1pt);
    \node[below left,fill=white,inner sep=1pt] at (1.9,0) {$A$};
    \fill (6.6,0) circle (1.1pt);
    \node[above,fill=white,inner sep=1pt] at (6.6,0.05) {$P$};
    \draw[thickline] (2.72,2.25) -- (7.0,-1.34);
    \draw[thickline] (2.72,-2.25) -- (7.0,1.34);
    \node at (4.3,-2.6) {(2)};
  \end{scope}
\end{tikzpicture}
\end{figbox}
\end{solution}

\begin{solution}{4.14}
(1) If the amplitudes of the incident, reflected and transmitted waves are as shown in
Fig. 4.14(a)(1) (the amplitude of the incident wave has already been set equal to 1), then
time-reversal invariance requires that the situation of Fig. 4.14(a)(2) is also possible. The
superposition principle, on the other hand, allows the situation of Fig. 4.14(a)(2) to be
regarded as the sum of Fig. 4.14(a)(3) and Fig. 4.14(a)(4). Making the incident and emergent
waves of this sum equal to the incident and emergent waves of Fig. 4.14(a)(2) gives the Stokes
relations
\begin{align*}
r^2 + tt' &= 1,\\
r' + r &= 0.
\end{align*}

(2) Let $A$ denote the amplitude of the incident wave I, and let $r_1$, $r_1'$, $t$, $t'$
denote respectively the (amplitude) reflection and transmission coefficients at the first
interface, while $r_2$ denotes the reflection coefficient at the second interface. Then the
amplitudes of the reflected waves $R_0, R_1, R_2, R_3, \ldots$ (Fig. 4.14(b); for clarity the
incidence is drawn obliquely) are as given in the following table:

\begin{center}
\renewcommand{\arraystretch}{1.35}
\begin{tabular}{cccccc}
$R_0$ & $R_1$ & $R_2$ & $R_3$ & $R_4$ & $\cdots$\\ \hline
$Ar_1$ & $Att'r_2$ & $Att'r_2r_1'r_2$ & $Att'r_2(r_1'r_2)^2$ & $Att'r_2(r_1'r_2)^3$ & $\cdots$
\end{tabular}
\end{center}

Taking into account the phase changes of the various reflected waves, the total amplitude of the
wave reflected into vacuum is
\begin{align*}
A_R &= Ar_1 + Att'r_2\ee^{-2\ii kd}\left[1 + r_1'r_2\ee^{-2\ii kd}
      + (r_1'r_2)^2\ee^{-4\ii kd} + \cdots\right]\\
    &= Ar_1 + \frac{Att'r_2\ee^{-2\ii kd}}{1 - r_1'r_2\ee^{-2\ii kd}}\\
    &= \frac{A\left(r_1 - r_1r_1'r_2\ee^{-2\ii kd} + tt'r_2\ee^{-2\ii kd}\right)}
            {1 - r_1'r_2\ee^{-2\ii kd}}\\
    &= \frac{A\left[r_1 + r_2\ee^{-2\ii kd}(tt' - r_1r_1')\right]}
            {1 - r_1'r_2\ee^{-2\ii kd}} .
\end{align*}

\begin{figbox}{4.14(a)}
\begin{tikzpicture}[scale=0.92]
  \foreach \dx/\lab in {0/(1), 3.0/(2), 6.0/(3), 9.0/(4)}{
    \begin{scope}[xshift=\dx cm]
      \draw[thickline] (-1.0,0) -- (1.0,0);
      \node[left] at (-1.0,0.30) {1};
      \node[left] at (-1.0,-0.30) {2};
      \node at (0,-1.75) {\lab};
    \end{scope}
  }
  % panel (1)
  \begin{scope}
    \draw[force] (0,1.15) -- (0,0.05);
    \node[above] at (0,1.18) {$1$};
    \draw[force] (0.55,0.05) -- (0.55,1.15);
    \node[above] at (0.55,1.18) {$r$};
    \draw[force] (-0.55,-0.05) -- (-0.55,-1.15);
    \node[below] at (-0.55,-1.18) {$t$};
  \end{scope}
  % panel (2)
  \begin{scope}[xshift=3.0cm]
    \draw[force] (0.55,0.05) -- (0.55,1.15);
    \node[above] at (0.55,1.18) {$1$};
    \draw[force] (0,1.15) -- (0,0.05);
    \node[above] at (0,1.18) {$r$};
    \draw[force] (-0.55,-0.05) -- (-0.55,-1.15);
    \node[below] at (-0.55,-1.18) {$t$};
  \end{scope}
  % panel (3)
  \begin{scope}[xshift=6.0cm]
    \draw[force] (0,-1.15) -- (0,-0.05);
    \node[below] at (0,-1.18) {$t$};
    \draw[force] (0.55,0.05) -- (0.55,1.15);
    \node[above] at (0.55,1.18) {$tt'$};
    \draw[force] (-0.55,-0.05) -- (-0.55,-1.15);
    \node[below] at (-0.55,-1.18) {$tr'$};
  \end{scope}
  % panel (4)
  \begin{scope}[xshift=9.0cm]
    \draw[force] (0,1.15) -- (0,0.05);
    \node[above] at (0,1.18) {$r$};
    \draw[force] (0.55,0.05) -- (0.55,1.15);
    \node[above] at (0.55,1.18) {$r^2$};
    \draw[force] (-0.55,-0.05) -- (-0.55,-1.15);
    \node[below] at (-0.55,-1.18) {$rt$};
  \end{scope}
\end{tikzpicture}
\end{figbox}

\begin{figbox}{4.14(b)}
\begin{tikzpicture}[scale=0.9]
  % media
  \draw[thickline] (0,1.5) -- (9.9,1.5);
  \draw[thickline] (0,0.5) -- (9.9,0.5);
  \node at (9.1,1.72) {$n=1$};
  \node at (0.8,1.0) {$n_1$};
  \node at (9.2,0.05) {$n_2$};
  % thickness
  \draw[dim] (-0.45,0.5) -- (-0.45,1.5) node[midway,left] {$d$};
  % incident ray
  \draw[force] (0.7,3.05) -- (2.0,1.5);
  \node[left] at (0.85,2.85) {I};
  % multiple reflections
  \draw[force] (2.0,1.5) -- (3.05,3.05);
  \node[above] at (2.52,2.86) {$R_0$};
  \draw[thickline] (2.0,1.5) -- (2.8,0.5);
  \draw[thickline] (2.8,0.5) -- (3.6,1.5);
  \draw[force] (3.6,1.5) -- (4.65,3.05);
  \node[above] at (4.12,2.86) {$R_1$};
  \draw[thickline] (3.6,1.5) -- (4.4,0.5);
  \draw[thickline] (4.4,0.5) -- (5.2,1.5);
  \draw[force] (5.2,1.5) -- (6.25,3.05);
  \node[above] at (5.72,2.86) {$R_2$};
  \draw[thickline] (5.2,1.5) -- (6.0,0.5);
  \draw[thickline] (6.0,0.5) -- (6.8,1.5);
  \draw[force] (6.8,1.5) -- (7.85,3.05);
  \node[above] at (7.32,2.86) {$R_3$};
  \draw[thickline] (6.8,1.5) -- (7.6,0.5);
  \draw[thickline] (7.6,0.5) -- (8.4,1.5);
  \draw[force] (8.4,1.5) -- (9.45,3.05);
  \node[above] at (8.92,2.86) {$R_4$};
  \node at (9.55,2.86) {$\cdots$};
\end{tikzpicture}
\end{figbox}

Here $k = 2\pi n_1/\lambda$, $\lambda$ being the wavelength in vacuum. The factor
$\ee^{-2\ii kd}$ is the phase change of the wave after it has traversed a distance $2d$ in the
medium of refractive index $n_1$. Using the Stokes relations derived in (1), we obtain
\begin{equation*}
\frac{A_R}{A} = \frac{r_1 + r_2\ee^{-2\ii kd}}{1 + r_1r_2\ee^{-2\ii kd}} .
\end{equation*}
From the above expression, the condition for the reflected wave to vanish is
\begin{equation*}
\text{1)}\quad
\begin{cases}
r_1 = -r_2,\\
\ee^{-2\ii kd} = 1,
\end{cases}
\qquad\text{or}\qquad
\text{2)}\quad
\begin{cases}
r_1 = r_2,\\
\ee^{-2\ii kd} = -1.
\end{cases}
\end{equation*}
Let us connect the coefficients $r_1$, $r_2$ with the refractive indices. By the Fresnel
formulae,
\begin{equation*}
r_1 = \frac{1-n_1}{1+n_1},\qquad
r_2 = \frac{n_1-n_2}{n_1+n_2}.
\end{equation*}
Expressed in terms of the refractive indices, the condition for the reflected wave to vanish
becomes
\begin{equation*}
\text{1)}\quad
\begin{cases}
n_2 = 1,\\[2pt]
\dfrac{n_1d}{\lambda} = \dfrac{p}{2},
\end{cases}
\qquad\text{or}\qquad
\text{2)}\quad
\begin{cases}
n_2 = n_1^2,\\[2pt]
\dfrac{n_1d}{\lambda} = \dfrac{2p-1}{4},
\end{cases}
\end{equation*}
where $p$ is a positive integer. The physical meaning of the above conditions is evident.

\begin{commentary}
This problem is precisely the physical principle of the antireflection coating. The result agrees
with that obtained by considering only the interference of the two reflected waves from the upper
and lower surfaces. This is because at normal incidence the reflection coefficients are very
small, so that apart from these first two reflected waves, which are relatively strong, all the
remaining reflected waves are very weak.
\end{commentary}
\end{solution}

\begin{solution}{4.15}
The speed of light in the lower tube is $c/n$, while its speed in the upper tube (relative to
the laboratory frame) is given by the relativistic velocity-addition formula as
\begin{equation*}
v' = \frac{\dfrac{c}{n}+u}{1+\dfrac{1}{c^2}\dfrac{c}{n}u}
   = \frac{\dfrac{c}{n}+u}{1+\dfrac{u}{cn}} .
\end{equation*}
Hence the equivalent refractive index of the upper tube is
\begin{align*}
n' &= \frac{c}{v'} = \frac{c}{\dfrac{c}{n}+u}
     \left(1+\frac{u}{cn}\right)
   = \frac{n}{1+\dfrac{nu}{c}}\left(1+\frac{u}{cn}\right)\\
   &\approx n\left(1+\frac{u}{cn}\right)\left(1-\frac{u}{cn}\right)
   \approx n-\frac{u}{c}(n^2-1).
\end{align*}
Therefore the optical path difference of the light in the upper and lower tubes is
\begin{equation*}
\Delta = L(n-n') = \frac{u}{c}(n^2-1)L .
\end{equation*}
When $\Delta = (2p-1)\frac{\lambda}{2}$ ($p$ a positive integer), destructive interference
occurs at $S'$. Setting $p=1$, we obtain the minimum value of $u$:
\begin{equation*}
u_{\min} = \frac{1}{2}\frac{\lambda}{L}\frac{c}{n^2-1} .
\end{equation*}
\end{solution}

\begin{solution}{4.16}
Let $\delta = \frac{2\pi}{\lambda}d\sin\theta$ be the phase difference between the light rays
emitted from two adjacent slits in the direction $\theta$, and let $a$ be the amplitude coming
from a single slit in that direction. Then the total amplitude is
\begin{equation*}
A = a\left[(1+\alpha) + (1-\alpha)\ee^{\ii\delta} + (1+\alpha)\ee^{\ii 2\delta}
    + (1-\alpha)\ee^{\ii 3\delta} + \cdots + (1-\alpha)\ee^{\ii(N-1)\delta}\right].
\end{equation*}
Here we take $N$ to be even. After rearranging, we have
\begin{align*}
A &= a\left[\sum_{n=0}^{N-1}\ee^{\ii n\delta}
     + \alpha\sum_{n=0}^{N-1}(-1)^n\ee^{\ii n\delta}\right]\\
  &= a\left[\frac{1-\ee^{\ii N\delta}}{1-\ee^{\ii\delta}}
     + \alpha\frac{1-\ee^{\ii N\delta}}{1+\ee^{\ii\delta}}\right]\\
  &= a\,\frac{\ee^{\ii N\delta/2}}{\ee^{\ii\delta/2}}
     \left[\frac{\ee^{-\ii N\delta/2}-\ee^{\ii N\delta/2}}
                {\ee^{-\ii\delta/2}-\ee^{\ii\delta/2}}
     + \alpha\frac{\ee^{-\ii N\delta/2}-\ee^{\ii N\delta/2}}
                {\ee^{-\ii\delta/2}+\ee^{\ii\delta/2}}\right]\\
  &= a\ee^{\ii(N-1)\delta/2}
     \left[\frac{\sin(N\delta/2)}{\sin(\delta/2)}
     - \ii\alpha\frac{\sin(N\delta/2)}{\cos(\delta/2)}\right].
\end{align*}
The intensity is proportional to
\begin{equation*}
AA^* = a^2\left[\frac{\sin^2(N\delta/2)}{\sin^2(\delta/2)}
      + \alpha^2\frac{\sin^2(N\delta/2)}{\cos^2(\delta/2)}\right].
\end{equation*}
Noting that $I_0\propto N^2a^2$ is the central maximum, the above expression can be written as
\begin{equation*}
I = I_0\,\frac{\sin^2(N\delta/2)}{N^2\sin^2(\delta/2)}
    \left[1+\alpha^2\tan^2\frac{\delta}{2}\right].
\end{equation*}
When $\delta \approx (2p+1)\pi$ ($p$ an integer), the effect of $\alpha$ is most easily detected.
\end{solution}

\begin{solution}{4.17}
The displacement of the $i$-th step of the drunkard can be represented by the vector
$\vec{L}_i$. After a certain time, the vector $\vec{R}$ from the hotel to the drunkard is the sum
of many vectors $\vec{L}_i$ of various directions:
\begin{equation*}
\vec{R} = \sum_{i=1}^{N}\vec{L}_i .
\end{equation*}
$N$ is the number of steps taken by the drunkard during this time and is a large number. What we
require is actually the modulus of the vector $\vec{R}$. Taking the hotel as the origin and
choosing $x$, $y$ coordinates arbitrarily, we have
\begin{equation*}
\vec{L}_i = L\cos\alpha_i\,\hat{x} + L\sin\alpha_i\,\hat{y},
\end{equation*}
where $\alpha_i$ is the angle between $\vec{L}_i$ and the $x$-axis (as shown in the figure).
Then
\begin{align*}
R^2 = R_x^2+R_y^2
&= \left(\sum_{i=1}^{N}L\cos\alpha_i\right)^2
 + \left(\sum_{i=1}^{N}L\sin\alpha_i\right)^2\\
&= \sum_{i=1}^{N}L^2\cos^2\alpha_i + \sum_{i\neq j}L^2\cos\alpha_i\cos\alpha_j
 + \sum_{i=1}^{N}L^2\sin^2\alpha_i + \sum_{i\neq j}L^2\sin\alpha_i\sin\alpha_j\\
&= L^2\sum_{i=1}^{N}\left(\cos^2\alpha_i+\sin^2\alpha_i\right)
 + L^2\sum_{i\neq j}\left(\cos\alpha_i\cos\alpha_j+\sin\alpha_i\sin\alpha_j\right).
\end{align*}
Since the $\alpha_i$ are arbitrary and $N$ is very large, the second term above evidently
vanishes through the cancellation of positive and negative contributions, while the first term
is evidently $NL^2$.

\begin{figbox}{4.17}
\begin{tikzpicture}[scale=1.15]
  \draw[axis] (-0.3,0) -- (6.6,0) node[right] {$x$};
  \draw[axis] (0,-0.3) -- (0,4.4) node[above] {$y$};
  \node[below left] at (0,0) {$O$};
  % random path
  \coordinate (P0) at (0,0);
  \coordinate (P1) at (1.7,0.55);
  \coordinate (P2) at (2.6,1.55);
  \coordinate (P3) at (4.3,1.35);
  \coordinate (P4) at (5.2,2.55);
  \coordinate (P5) at (4.5,3.75);
  \draw[dashedline] (P0) -- (P1) -- (P2) -- (P3) -- (P4) -- (P5);
  \draw[force] (P0) -- (P1);
  \draw[force] (P1) -- (P2);
  \draw[force] (P2) -- (P3);
  \draw[force] (P3) -- (P4);
  \draw[force] (P4) -- (P5);
  % labels
  \node[fill=white,inner sep=1pt] at (0.85,0.44) {$L_1$};
  \node[fill=white,inner sep=1pt] at (2.42,1.32) {$L_2$};
  \node[fill=white,inner sep=1pt] at (3.45,1.10) {$L_3$};
  \node[fill=white,inner sep=1pt] at (5.02,1.78) {$L_4$};
  \node[fill=white,inner sep=1pt] at (4.28,3.38) {$L_5$};
  % resultant
  \draw[force,dashed] (P0) -- (P5);
  \node[fill=white,inner sep=1pt] at (2.05,1.95) {$R$};
  % angles
  \draw (0.55,0) arc (0:17.9:0.55);
  \node at (0.80,0.16) {$\alpha_1$};
  \draw (P1) ++(17.9:0.45) arc (17.9:48.0:0.45);
  \node at (2.02,0.78) {$\alpha_2$};
  \draw (P2) ++(48.0:0.45) arc (48.0:-6.7:0.45);
  \node at (2.95,1.80) {$\alpha_3$};
  \draw (P3) ++(-6.7:0.45) arc (-6.7:53.1:0.45);
  \node at (4.60,1.78) {$\alpha_4$};
  \draw (P4) ++(53.1:0.45) arc (53.1:120.3:0.45);
  \node at (4.88,2.90) {$\alpha_5$};
\end{tikzpicture}
\end{figbox}

Hence
\begin{equation*}
R^2 = NL^2,
\end{equation*}
so that
\begin{equation*}
R = \sqrt{N}L .
\end{equation*}
If the drunkard takes the same time for every step, then $N\propto t$, and hence
\begin{equation*}
R \propto \sqrt{t}\,L,
\end{equation*}
that is, the distance between the drunkard and the hotel is proportional to the square root of
the time.

\begin{commentary}
The physical ideas and methods of this problem are manifested in many fields of physics. In
optics it corresponds to the superposition of $N$ light waves whose phases vary at random, the
result being the absence of any interference phenomenon; in molecular physics it is the physical
basis of Brownian motion.
\end{commentary}
\end{solution}

\begin{solution}{4.18}
Let $O$ be the geometrical focus of the rays (as shown in the figure). Let $M$ be the plane
through $O$ perpendicular to the principal axis of the lens, and let $r_0$ be the shortest
distance from the lens to the point $O$. Draw a spherical surface $S$ with centre $O$ and radius
$r_0$, its edge being defined by the aperture stop. When calculating the light field on the
surface $S$, geometrical optics may be applied approximately. In this way the light field on $S$
can be written in complex form as
\begin{equation*}
E_S = \frac{1}{r_0}\ee^{\ii(\omega t+k r_0)} .
\end{equation*}
Let $P$ be an observation point on the plane $M$, $\dd S$ an area element on $S$, and $r$ the
distance between $\dd S$ and $P$. According to Huygens' principle, the light field at the point
$P$ can be obtained from the following formula:
\begin{equation}
E_P = \int_S \frac{1}{r_0r}\ee^{\ii[\omega t-k(r-r_0)]}\,\dd S . \tag{1}
\end{equation}
Putting $OP=\rho$, we have $\vec{r}=\vec{r}_0+\vec{\rho}$, so that
$r^2 = r_0^2+\rho^2+2\vec{\rho}\cdot\vec{r}_0$ and
$r = r_0\sqrt{1+\frac{\rho^2}{r_0^2}+\frac{2\vec{\rho}\cdot\vec{r}_0}{r_0^2}}$. Since $\rho$ is
a small quantity, Newton's binomial formula is applicable, and to accuracy of order $\rho^2$,
\begin{equation}
r-r_0 = \vec{\rho}\cdot\hat{r}_0 + \frac{\rho^2}{2r_0}
       - \frac{1}{2}\frac{(\vec{\rho}\cdot\vec{r}_0)^2}{r_0^3} . \tag{2}
\end{equation}
When $\rho \ll r,\,r_0$, to first approximation,
\begin{equation*}
r-r_0 = \vec{\rho}\cdot\hat{r}_0 ,
\end{equation*}
where $\hat{r}_0$ is the unit vector in the direction of $\vec{r}_0$. The error permitted in
calculating $(r-r_0)$ must be much smaller than the wavelength; from (2),
$r-r_0-\vec{\rho}\cdot\vec{r}_0\approx \rho^2/r_0$, so when calculating $(r-r_0)$ the following
condition must be satisfied:
\begin{equation}
\rho^2/r_0 \ll \lambda . \tag{3}
\end{equation}
We now calculate the light field in the neighbourhood of $O$ under condition (3). Introduce
rectangular coordinates in the observation plane, with the origin at $O$, the coordinate axes
$x$ and $y$ parallel to adjacent sides of the square aperture, and the optical axis as the
$z$-axis. Let $x$ and $y$ denote the coordinates of $P$. The direction of the unit vector
$\hat{r}_0$ can be expressed by the angles $\varphi$ and $\psi$, where $\varphi$ and $\psi$ are
respectively the angles made by $r_0$ with the $yz$-plane and the $xz$-plane. Thus
\begin{equation*}
r-r_0 = \vec{\rho}\cdot\hat{r}_0 = x\sin\varphi + y\sin\psi \approx x\varphi + y\psi ,
\end{equation*}
\begin{equation*}
\dd S = r_0^2\,\dd\varphi\,\dd\psi .
\end{equation*}

\begin{figbox}{4.18}
\begin{tikzpicture}[scale=1.0]
  % lens
  \draw[thickline,fill=black!10]
    (1.0,2.2) .. controls (0.45,0.8) and (0.45,-0.8) .. (1.0,-2.2)
    .. controls (1.55,-0.8) and (1.55,0.8) .. cycle;
  % wavefront S
  \draw[thickline] (7,0) ++(152:6.0) arc (152:208:6.0);
  \node[right] at (1.62,-1.95) {$S$};
  \node[above left] at (1.40,2.33) {$\dd S$};
  \fill (1.44,2.25) circle (1.2pt);
  % plane M
  \draw[thickline] (7,-2.7) -- (7,2.8);
  \node[left] at (6.82,2.85) {$M$};
  \fill (7,0) circle (1.2pt);
  \node[below right] at (7,-0.02) {$O$};
  \fill (7,1.35) circle (1.1pt);
  \node[right] at (7,1.35) {$P$};
  % rays
  \draw[thickline] (7,0) -- (1.44,2.25) node[midway,above] {$r_0$};
  \draw[thickline] (1.44,2.25) -- (7,1.35) node[midway,above right] {$r$};
  \draw[dim] (7,0) -- (7,1.35) node[midway,right] {$\rho$};
  \draw[thinline,dashed] (1.0,2.2) -- (7,0);
  \draw[thinline,dashed] (1.0,-2.2) -- (7,0);
  % axes
  \draw[axis] (7,0) -- (0.3,0) node[left] {$z$};
  \draw[axis] (7,0) -- (8.9,-0.55) node[right] {$x$};
  \draw[axis] (7,0) -- (7,2.30) node[above right,fill=white,inner sep=1pt] {$y$};
\end{tikzpicture}
\end{figbox}

We assume that the half-angle subtended at $O$ by the square aperture,
$\alpha = \arctan\frac{D}{2r_0}$, is very small, so that $\sin\varphi$ and $\sin\psi$ may be
replaced by $\varphi$ and $\psi$. Apart from this, since $\rho \ll r,\,r_0$, the quantity $r$ in
the denominator of (1) may be replaced approximately by $r_0$. Finally we obtain
\begin{equation*}
E_P = \ee^{\ii\omega t}\int_{-\alpha}^{+\alpha}\!\!\int_{-\alpha}^{+\alpha}
      \ee^{-\ii k(x\varphi+y\psi)}\,\dd\varphi\,\dd\psi .
\end{equation*}
Integrating the above expression, taking the amplitude at $O$ as unity and noting that
$k = \frac{2\pi}{\lambda}$, we obtain the amplitude at $P$:
\begin{equation*}
A_P = \frac{\sin\dfrac{2\pi x\alpha}{\lambda}}{\dfrac{2\pi x\alpha}{\lambda}}
      \cdot
      \frac{\sin\dfrac{2\pi y\alpha}{\lambda}}{\dfrac{2\pi y\alpha}{\lambda}} ,
\end{equation*}
and the intensity is
\begin{equation*}
I_P = \left[\frac{\sin\dfrac{2\pi x\alpha}{\lambda}}
             {\dfrac{2\pi x\alpha}{\lambda}}\right]^2
      \left[\frac{\sin\dfrac{2\pi y\alpha}{\lambda}}
             {\dfrac{2\pi y\alpha}{\lambda}}\right]^2 .
\end{equation*}
The diffraction pattern obtained is the same as the Fraunhofer diffraction pattern of a square
aperture. The distance from the centre of the central maximum to the first minimum, and the
distance between two adjacent minima, are both equal to
\begin{equation*}
\Delta x = \Delta y = \frac{\lambda/2}{\alpha} .
\end{equation*}
It remains to check whether condition (3) is satisfied. Since the range of perceptible intensity
is $\Delta x$, we may take $\rho \approx \lambda/\alpha$; substituting into (3) changes the
condition into
\begin{equation*}
\alpha \gg \sqrt{\lambda/r_0} .
\end{equation*}
This condition is satisfied in all optical instruments constructed from lenses or plane mirrors.

\begin{commentary}
This problem concerns the basic method for calculating a diffracted light field. The result shows
that it is not necessary to satisfy strictly the Fraunhofer diffraction conditions (the source
and the image field both at infinity, or both in the focal plane of a lens) in order to observe
a Fraunhofer diffraction pattern; it can also be observed in the geometrical image plane of the
source, provided that the geometrical conditions for imaging are satisfied and the observation
region lies near the geometrical image point.
\end{commentary}
\end{solution}

\begin{solution}{4.19}
(1) Let the intensities of the elliptically polarised component along the $x$- and $y$-axes be
$I_{ex}$ and $I_{ey}$ respectively; the unpolarised component has the same intensity along the
$x$- and $y$-axes, which we denote by $I_u$. Then from the conditions of the problem,
\begin{align}
I_u + I_{ex} &= 1.5I_0, \tag{1}\\
I_u + I_{ey} &= 1.0I_0. \tag{2}
\end{align}
When the transmission axis of the polariser makes an angle $\theta$ with the $x$-axis, the
unpolarised component passing through the polariser is still $I_u$, while the amplitude of the
elliptically polarised component passing through it is
\begin{equation*}
E_\theta = E_x\cos\theta + E_y\sin\theta .
\end{equation*}
Since $x$ and $y$ are the principal axes of the ellipse, the phase difference between $E_x$ and
$E_y$ is $90^\circ$, and hence
$E_\theta^2 = E_x^2\cos^2\theta + E_y^2\sin^2\theta$; therefore
\begin{align*}
I_\theta &= I_{ex}\cos^2\theta + I_{ey}\sin^2\theta + I_u\\
          &= (I_u+I_{ex})\cos^2\theta + (I_u+I_{ey})\sin^2\theta\\
          &= 1.5I_0\cos^2\theta + 1.0I_0\sin^2\theta .
\end{align*}
The result is independent of the unpolarised component.

(2) After passing through the quarter-wave plate, the phase difference between $E_x$ and $E_y$ is
changed by $90^\circ$, so that they become in phase. Since the transmitted intensity is maximum
when the polariser makes an angle of $30^\circ$ with the $x$-axis, it is seen that
\begin{equation*}
\frac{E_x}{E_y} = \cot 30^\circ = \sqrt{3},
\qquad\text{i.e.}\qquad
\frac{I_{ex}}{I_{ey}} = 3 .
\end{equation*}
Substituting into (1), we obtain
\begin{align*}
I_{ex} &= 0.75I_0,\\
I_{ey} &= 0.25I_0,\\
I_u &= 0.75I_0 .
\end{align*}
Hence the maximum intensity at $\theta = 30^\circ$ is
\begin{equation*}
I = I_{ex}+I_{ey}+I_u = 1.75I_0 .
\end{equation*}
The percentage of the unpolarised component in the total incident intensity is
\begin{equation*}
\frac{2I_u}{2I_u+I_{ex}+I_{ey}}
= \frac{2\times 0.75I_0}{2\times 0.75I_0+0.75I_0+0.25I_0} = 60\% .
\end{equation*}
\end{solution}

\begin{solution}{4.20}
In a birefringent crystal, light waves polarised parallel and perpendicular to the optic axis
travel with different phase velocities, so that the linearly polarised wave
$E_0(\hat{x}+\hat{y})/\sqrt{2}$ emerging from the polarising film becomes, after entering the
crystal and travelling a distance $z$,
\begin{equation*}
\vec{E} = \frac{E_0}{\sqrt{2}}
          \left(\hat{x}\,\ee^{\ii k n_+ z} + \hat{y}\,\ee^{\ii k n_- z}\right)\ee^{-\ii\omega t},
\end{equation*}
where $\hat{y}$ is chosen parallel to the optic axis, $\hat{x}$ perpendicular to $\hat{y}$ and
to the direction of propagation $\hat{z}$, the phase velocity of the wave polarised along $x$ is
$c/n_+$, and that of the wave polarised along $y$ is $c/n_-$; the plane of the polariser is
parallel to $\vec{l} = (\hat{x}+\hat{y})/\sqrt{2}$. After travelling the whole length $2^N d$ of
the crystal, the amplitude of the wave transmitted through the polariser is $\vec{E}\cdot\vec{l}$,
namely
\begin{equation*}
\vec{E}\cdot\vec{l}
= \frac{E_0}{2}\left(\ee^{\ii k n_+ 2^N d} + \ee^{\ii k n_- 2^N d}\right).
\end{equation*}
With
\begin{equation*}
n_+ = \frac{1}{2}\left[(n_+ + n_-) + (n_+ - n_-)\right],\qquad
n_- = \frac{1}{2}\left[(n_+ + n_-) - (n_+ - n_-)\right],
\end{equation*}
substitution into the above gives
\begin{align*}
\vec{E}\cdot\vec{l}
&= \frac{E_0}{2}\left\{\ee^{\ii k 2^{N-1}d\left[(n_+ + n_-)+(n_+ - n_-)\right]}
   + \ee^{\ii k 2^{N-1}d\left[(n_+ + n_-)-(n_+ - n_-)\right]}\right\}\\
&= \frac{E_0}{2}\,\ee^{\ii k 2^{N-1}d(n_+ + n_-)}
   \left[\ee^{\ii k 2^{N-1}d(n_+ - n_-)}
   + \ee^{-\ii k 2^{N-1}d(n_+ - n_-)}\right]\\
&= \frac{E_0}{2}\,\ee^{\ii k 2^{N-1}d(n_+ + n_-)}
   \left[\ee^{\ii k 2^N d(n_+ - n_-)/2}
   + \ee^{-\ii k 2^N d(n_+ - n_-)/2}\right]\\
&= E_0\,\ee^{\ii k 2^{N-1}d(n_+ + n_-)}\cos\left(2^N\varphi\right).
\end{align*}
\end{solution}

